\chapter{Dyspnoea}
\index{Dyspnoea}

\section*{Definition and Overview}
Dyspnoea is the medical term for breathlessness, the unpleasant sensation of not being able to breathe easily or enough. It is a common symptom that ranges from mild breathlessness on exertion to severe difficulty breathing at rest. Dyspnoea is a symptom of an underlying problem rather than a disease in itself, and the causes include heart and lung conditions, anaemia, anxiety, and physical deconditioning. Because breathlessness can signal a serious or life-threatening condition, it always warrants careful assessment, and sudden or severe breathlessness is a medical emergency.

\section*{Epidemiology}
Breathlessness is one of the most common reasons for seeking medical care and is a leading symptom in emergency departments. It becomes more common with age, and it affects a large proportion of people with chronic conditions such as chronic obstructive pulmonary disease, asthma, heart failure, and obesity. Breathlessness is also a major cause of disability and reduced quality of life, and it is associated with increased hospital admission and mortality.

\section*{Aetiology and Pathophysiology}
Breathlessness arises when the body's demand for oxygen exceeds its ability to supply it, or when the brain perceives that breathing is inadequate.
\begin{itemize}
    \item \textbf{Lung conditions:} asthma, chronic obstructive pulmonary disease, pneumonia, pulmonary embolism, and lung cancer
    \item \textbf{Heart conditions:} heart failure, in which the heart cannot pump effectively, and coronary artery disease
    \item \textbf{Anaemia:} reduced oxygen-carrying capacity of the blood
    \item \textbf{Metabolic and other causes:} thyroid disease, kidney failure, and severe obesity
    \item \textbf{Anxiety and panic:} hyperventilation and the sensation of air hunger
    \item \textbf{Deconditioning:} poor physical fitness makes normal exertion feel breathless
    \item \textbf{The sensation of dyspnoea:} results from the interplay of chemical signals in the blood, lung stretch receptors, and the brain's perception of effort, so the severity of breathlessness does not always match the severity of the underlying disease
\end{itemize}

\section*{Clinical Features}
The pattern of breathlessness provides important clues to the cause.
\begin{itemize}
    \item Breathlessness on exertion that limits activity
    \item Breathlessness at rest or after minimal activity
    \item Sudden onset of severe breathlessness, which may indicate a clot or collapsed lung
    \item Wheezing, cough, and sputum production in lung disease
    \item Swelling of the ankles, orthopnoea, and waking breathless at night in heart failure
    \item Chest pain, palpitations, or faintness with breathlessness
    \item Pale skin and fatigue in anaemia
    \item Anxiety, hyperventilation, and a feeling of needing to take deep breaths
\end{itemize}

\section*{Differential Diagnosis}
The many causes of breathlessness must be systematically considered.
\begin{itemize}
    \item Asthma and chronic obstructive pulmonary disease
    \item Heart failure and cardiac ischaemia
    \item Pulmonary embolism, particularly with sudden onset and risk factors
    \item Pneumonia and other respiratory infections
    \item Anaemia and thyroid disorders
    \item Anxiety and panic disorder
    \item Obesity and deconditioning
    \item Rare causes including pulmonary fibrosis and pulmonary hypertension
\end{itemize}

\section*{When to Seek Medical Care}
New, persistent, or worsening breathlessness always requires medical assessment, even if the cause turns out to be minor. People with known lung or heart conditions should follow their action plans and seek review when symptoms change.

\textbf{Contact your local emergency services for:}
\begin{itemize}
    \item Sudden severe breathlessness or breathlessness at rest
    \item Breathlessness with chest pain, which may indicate a heart attack or pulmonary embolism
    \item Breathlessness with bluish lips or fingers
    \item Difficulty speaking full sentences because of breathlessness
    \item Collapse, confusion, or severe faintness
\end{itemize}

\section*{Diagnosis and Investigations}
The investigation of breathlessness is guided by the history and examination.
\begin{itemize}
    \item \textbf{History:} onset, timing, severity, and associated symptoms
    \item \textbf{Examination:} oxygen saturation, respiratory rate, lung sounds, heart rate, and signs of fluid overload
    \item \textbf{Spirometry:} lung function testing for asthma and COPD
    \item \textbf{Chest X-ray:} to identify infection, fluid, lung collapse, or masses
    \item \textbf{Blood tests:} full blood count for anaemia, and tests for heart failure
    \item \textbf{ECG and echocardiogram:} to assess the heart
    \item \textbf{Further imaging:} CT scans for suspected pulmonary embolism or other diagnoses
\end{itemize}

\section*{Management}
Management addresses the underlying cause and relieves symptoms.
\begin{itemize}
    \item \textbf{Treatment of the cause:} inhalers for asthma and COPD, diuretics for heart failure, antibiotics for pneumonia, blood thinners for pulmonary embolism, and iron for anaemia
    \item \textbf{Oxygen therapy:} for people with low oxygen levels, including home oxygen in chronic conditions
    \item \textbf{Pulmonary rehabilitation:} exercise training and education that improve breathlessness in chronic lung disease
    \item \textbf{Breathing techniques:} positions and methods that reduce the work of breathing
    \item \textbf{Management of anxiety:} treatment of panic and anxiety, which can drive or worsen breathlessness
    \item \textbf{Lifestyle measures:} weight loss, smoking cessation, and gradual reconditioning
    \item \textbf{Palliative approaches:} symptom control in advanced disease
\end{itemize}

\section*{Complications}
\begin{itemize}
    \item Reduced activity and disability
    \item Falls and fractures from reduced mobility
    \item Anxiety and depression related to chronic breathlessness
    \item Social isolation
    \item Hospital admission and the complications of the underlying disease
    \item Missed diagnosis of serious conditions if symptoms are dismissed
\end{itemize}

\section*{Prognosis and Outcomes}
The outlook for dyspnoea depends on the cause. Breathlessness from treatable causes such as anaemia, pneumonia, and deconditioning resolves with treatment. Chronic breathlessness from heart or lung disease can be substantially improved with optimal management and rehabilitation, although it may not fully resolve. Early and accurate diagnosis is important, because conditions such as pulmonary embolism are life-threatening but treatable.

\section*{Prevention and Self-Care}
\begin{itemize}
    \item Do not smoke, and avoid exposure to smoke and pollutants
    \item Maintain physical activity and fitness appropriate to your health
    \item Take prescribed inhalers and other medicines as directed
    \item Follow a pulmonary rehabilitation or cardiac rehabilitation programme if recommended
    \item Use breathing techniques such as pursed-lip breathing during exertion
    \item Report any new or worsening breathlessness promptly
    \item Know the emergency signs and when to call for help
\end{itemize}

\section*{Sources and Further Reading}
The following reputable sources provide further information for healthcare students and patients seeking reliable, current advice about breathlessness.
\begin{itemize}
    \item Lung Foundation Australia. \textit{Breathlessness information and support.} https://lungfoundation.com.au/
    \item Healthdirect Australia. \textit{Breathlessness.} https://www.healthdirect.gov.au/breathlessness
    \item NHS. \textit{Breathlessness: causes and treatment.} https://www.nhs.uk/conditions/breathlessness/
    \item Mayo Clinic. \textit{Shortness of breath: symptoms and causes.} https://www.mayoclinic.org/symptoms/shortness-of-breath/
    \item MSD Manual. \textit{Dyspnea: professional overview.} https://www.msdmanuals.com/
    \item American Thoracic Society. \textit{Patient information on breathlessness.} https://www.thoracic.org/
\end{itemize}
